\chapter{KHẢO SÁT VÀ PHÂN TÍCH YÊU CẦU}

Từ bài toán đã nêu ở Chương 1, chương này cụ thể hóa yêu cầu của hệ thống. Nội dung được trình bày theo trình tự: khảo sát hiện trạng, xác định tác nhân, mô tả use case tổng quan, phân tích các quy trình nghiệp vụ chính, sau đó tổng hợp yêu cầu chức năng và phi chức năng.

\section{Khảo sát hiện trạng}

Quản lý dự án sinh viên khác với quản lý một danh sách bài tập ngắn hạn. Một dự án thường kéo dài nhiều tuần hoặc nhiều tháng, có nhiều thành viên, nhiều mốc kế hoạch, nhiều nhiệm vụ chính thức và nhiều lần nộp tiến độ. Giảng viên cần biết dự án nào cần phản hồi ngay, dự án nào đang quá hạn, nhóm nào thiếu minh chứng và bài nộp nào đã được trả lại để sửa. Sinh viên cần một nơi rõ ràng để xem công việc được giao, nộp kết quả và đọc phản hồi. Quản trị viên cần bảo đảm tài khoản, vai trò và trạng thái người dùng không làm sai lệch dữ liệu dự án.

Các công cụ rời rạc có thể hỗ trợ một phần công việc, nhưng không tạo được một quy trình có trạng thái rõ ràng. Email lưu được nội dung trao đổi nhưng khó lọc theo nhiệm vụ. Nhóm chat phản hồi nhanh nhưng thông tin dễ bị trôi. Bảng tính dễ chỉnh sửa nhưng thiếu phân quyền và nhật ký thay đổi. Thư mục file lưu được sản phẩm nhưng không thể hiện sản phẩm đó đã được duyệt hay chưa. Vì vậy, hệ thống cần gom các đối tượng nghiệp vụ chính vào một luồng dữ liệu có quan hệ rõ ràng.

Phần khảo sát này dựa trên phân tích quy trình thường gặp khi giảng viên theo dõi nhiều nhóm dự án bằng các công cụ phổ biến nêu trên. Trọng tâm không phải thay thế mọi công cụ giao tiếp, mà là xác định những dữ liệu nào cần được quản lý có trạng thái: dự án, nhiệm vụ, bài nộp, quyết định duyệt và minh chứng.

\tableintro{tab:current-tool-comparison}{so sánh các phương pháp theo dõi dự án hiện tại và rút ra yêu cầu nghiệp vụ tương ứng cho UniTrack}
\begin{table}[H]
\caption{So sánh phương pháp hiện tại và yêu cầu nghiệp vụ của UniTrack}
\label{tab:current-tool-comparison}
\centering
\begin{tblr}{
  width=\textwidth,
  colspec={X[0.85,l]X[1.05,l]X[1.45,l]X[1.45,l]},
  rows={valign=t},
  row{1}={font=\bfseries,bg=black!7},
  hlines={0.35pt,black!24},
}
Công cụ & Điểm mạnh & Hạn chế trong đồ án dự án & Yêu cầu rút ra cho UniTrack \\
Email & Có lịch sử trao đổi, quen thuộc & Khó theo dõi trạng thái từng nhiệm vụ, dễ lẫn phiên bản bài nộp & Bài nộp phải gắn trực tiếp với nhiệm vụ và dự án \\
Nhóm chat & Phản hồi nhanh & Không có cấu trúc nghiệp vụ, khó tìm lại quyết định duyệt bài cũ & Quyết định duyệt bài cần có lịch sử riêng và trạng thái cuối cùng \\
Bảng tính & Tổng hợp được danh sách & Dễ sai khi nhiều người sửa, thiếu kiểm tra quyền & Server và cơ sở dữ liệu phải giữ ràng buộc, không tin hoàn toàn dữ liệu từ giao diện \\
Thư mục lưu file & Lưu sản phẩm, ảnh, báo cáo & Không biết file thuộc bài nộp nào, có thể bị sửa sau khi đã duyệt & Minh chứng phải gắn với bài nộp và chuyển sang chỉ đọc sau khi duyệt \\
UniTrack & Dự án, nhiệm vụ, bài nộp, duyệt bài và minh chứng trong cùng hệ thống & Cần xây dựng ứng dụng Web và vận hành dữ liệu, session, file storage & Tập trung vào dự án, phân quyền và vòng đời rõ ràng \\
\end{tblr}
\end{table}

\section{Tác nhân hệ thống}

UniTrack có ba vai trò chính. Mỗi vai trò xác định màn hình được truy cập, phạm vi dữ liệu được truy vấn, quyền ghi dữ liệu và cách dashboard được sắp xếp.

\tableintro{tab:actors-responsibilities}{mô tả ba tác nhân chính và phạm vi trách nhiệm của từng tác nhân trong hệ thống}
\begin{table}[H]
\caption{Tác nhân và phạm vi trách nhiệm trong UniTrack}
\label{tab:actors-responsibilities}
\centering
\begin{tblr}{
  width=\textwidth,
  colspec={X[0.85,l]X[1.45,l]X[1.80,l]},
  rows={valign=t},
  row{1}={font=\bfseries,bg=black!7},
  hlines={0.35pt,black!24},
}
Tác nhân & Mục tiêu sử dụng & Phạm vi thao tác chính \\
Quản trị viên & Duy trì tài khoản và quan sát dữ liệu hệ thống & Tạo tài khoản, đổi vai trò/trạng thái, đặt mật khẩu, xem dashboard toàn cục, xem hoặc cập nhật dự án theo quy tắc hiện tại \\
Giảng viên & Giám sát dự án và phản hồi tiến độ sinh viên & Tạo thư mục/dự án, quản lý nhóm, tạo mốc kế hoạch/nhiệm vụ, duyệt bài nộp, quản lý tài nguyên và minh chứng theo vòng đời dự án \\
Sinh viên & Theo dõi công việc được giao và nộp tiến độ & Xem dự án tham gia, xem nhiệm vụ, nộp bài, tải minh chứng cho bài đang chờ duyệt và xem phản hồi \\
\end{tblr}
\end{table}

\tableintro{tab:role-permission-matrix}{tổng hợp quyền theo vai trò cùng các điều kiện nghiệp vụ bắt buộc khi thực hiện thao tác}
\begin{longtblr}[
  caption={Ma trận quyền theo vai trò và điều kiện nghiệp vụ},
  label={tab:role-permission-matrix},
]{
  width=\textwidth,
  colspec={X[1.2,l]X[0.55,c]X[0.65,c]X[0.65,c]X[2.2,l]},
  rowhead=1,
  rows={valign=t},
  row{1}={font=\bfseries,bg=black!7},
  hlines={0.35pt,black!24},
}
Nghiệp vụ & Quản trị & Giảng viên & Sinh viên & Điều kiện quan trọng \\
Quản lý tài khoản & Có & Không & Không & Phải giữ lại quản trị viên active cuối cùng; đặt mật khẩu phải thu hồi session cũ \\
Tạo thư mục/dự án & Có & Có & Không & Chủ sở hữu thư mục phải khớp giảng viên phụ trách dự án khi gán dự án vào thư mục \\
Quản lý thành viên & Có & Có & Không & Chỉ thêm sinh viên đang hoạt động; trạng thái dự án phải cho phép thao tác \\
Tạo/sửa nhiệm vụ & Có & Có & Không & Dự án phải active khi tạo mới; mốc kế hoạch và sinh viên được giao phải cùng dự án \\
Nộp bài & Không & Không & Có & Sinh viên phải được phân công; dự án active; không có bài chờ duyệt trùng nhiệm vụ \\
Duyệt bài & Có & Có & Không & Người duyệt phải quản lý dự án; bài nộp đang chờ duyệt; dự án chưa archived \\
Đọc/tải minh chứng & Có & Có & Có & Người dùng phải còn quyền xem dự án; minh chứng đã duyệt vẫn đọc được \\
\end{longtblr}

\section{Tổng quan chức năng}

\figurelead{fig:use-case-overview.png}{thể hiện phạm vi chức năng chính của UniTrack theo ba vai trò người dùng và các nhóm nghiệp vụ trọng tâm}
\figUseCaseOverview

\subsection{Nhóm quản lý tài khoản}

UniTrack không có đăng ký công khai. Quản trị viên tạo tài khoản để kiểm soát người dùng tham gia hệ thống. Cách tiếp cận này phù hợp với bối cảnh quản lý dự án trong trường học, nơi danh sách sinh viên/giảng viên cần được kiểm soát và tài khoản đã ngừng hoạt động phải bị thu hồi session cũ.

Nhóm chức năng này được tách thành hai phần: session dùng cho mọi người dùng và thao tác quản trị tài khoản dành riêng cho quản trị viên. Cách tách này giúp biểu đồ giữ đúng vai trò của từng tác nhân, đồng thời tránh đưa quá nhiều luật bảo vệ vào một hình.

\figurelead{fig:account-session-use-cases.png}{mô tả các tương tác xác thực cơ bản mà mọi người dùng đều phải đi qua trước khi sử dụng hệ thống}
\figAccountSessionUseCases

\figurelead{fig:admin-account-use-cases.png}{làm rõ các thao tác quản trị tài khoản và phạm vi quyền chỉ dành cho quản trị viên}
\figAdminAccountUseCases

\subsection{Nhóm Workspace và dự án}

Workspace là khu vực điều hướng chính sau dashboard. Giảng viên và quản trị viên nhìn thấy các thư mục lớp và danh sách dự án độc lập; sinh viên chỉ nhìn thấy dự án mình tham gia. Thư mục trong hệ thống được dùng như một cách nhóm dự án nhẹ, không mở rộng thành module khóa học độc lập.

Các use case của Workspace được chia thành phần tổ chức thư mục và phần quản lý dự án/nhóm. Phần thư mục tập trung vào việc xem, tìm kiếm và gán dự án vào đúng phạm vi giảng viên phụ trách; phần dự án/nhóm tập trung vào hồ sơ dự án, vòng đời và thành viên.

\figurelead{fig:workspace-folder-use-cases.png}{trình bày nhóm chức năng tổ chức Workspace, trong đó thư mục được sử dụng để gom các dự án theo phạm vi quản lý}
\figWorkspaceFolderUseCases

\figurelead{fig:project-team-use-cases.png}{thể hiện các thao tác quản lý hồ sơ dự án, vòng đời và nhóm thành viên trong phạm vi một dự án cụ thể}
\figProjectTeamUseCases

\subsection{Nhóm nhiệm vụ, bài nộp và duyệt bài}

Đây là quy trình lõi của sản phẩm. Nhiệm vụ phải thuộc mốc kế hoạch trong cùng dự án; sinh viên được giao phải là thành viên đang hoạt động của dự án; sinh viên chỉ được nộp khi dự án active và chưa có bài nộp nào đang chờ duyệt cho nhiệm vụ đó; quyết định duyệt bài là trạng thái cuối cùng của một bài nộp.

Để giữ sơ đồ rõ ràng, quy trình được tách thành hai nhóm use case. Nhóm thứ nhất mô tả lập kế hoạch và giao nhiệm vụ. Nhóm thứ hai mô tả bài nộp, tài nguyên/minh chứng và quyết định duyệt bài.

\figurelead{fig:assignment-plan-use-cases.png}{mô tả các chức năng lập kế hoạch, tạo mốc và giao nhiệm vụ chính thức cho sinh viên trong dự án}
\figAssignmentPlanUseCases

\figurelead{fig:submission-review-use-cases.png}{mô tả luồng sinh viên nộp bài, bổ sung minh chứng và giảng viên ghi nhận quyết định duyệt bài}
\figSubmissionReviewUseCases

\section{Danh sách use case và mô tả}

\tableintro{tab:use-case-list}{liệt kê các use case chính làm cơ sở chọn các use case tiêu biểu để đặc tả chi tiết}
\begin{longtblr}[
  caption={Danh sách use case chính của UniTrack},
  label={tab:use-case-list},
]{
  width=\textwidth,
  colspec={X[0.38,l]X[0.9,l]X[2.25,l]},
  rowhead=1,
  rows={valign=t},
  row{1}={font=\bfseries,bg=black!7},
  hlines={0.35pt,black!24},
}
ID & Tên use case & Mô tả \\
UC01 & Đăng nhập hệ thống & Người dùng có tài khoản active đăng nhập bằng email/mật khẩu để vào dashboard theo vai trò. \\
UC02 & Tạo dự án & Giảng viên/quản trị viên tạo dự án, chọn giảng viên phụ trách và gán thư mục nếu phù hợp. \\
UC03 & Tạo nhiệm vụ chính thức & Người quản lý tạo assignment thuộc mốc kế hoạch, đặt hạn nộp, ưu tiên và phân công sinh viên. \\
UC04 & Nộp báo cáo và minh chứng & Sinh viên được phân công gửi báo cáo tiến độ và gắn file minh chứng cho bài đang chờ duyệt. \\
UC05 & Duyệt bài nộp & Giảng viên/quản trị viên xem bài nộp, tài nguyên, minh chứng và ghi quyết định duyệt cuối cùng. \\
UC06 & Quản lý tài khoản & Quản trị viên tạo tài khoản, tìm kiếm, đổi vai trò/trạng thái và đặt mật khẩu. \\
UC07 & Thêm sinh viên vào dự án & Người quản lý thêm sinh viên active vào dự án và kiểm tra lại trạng thái tài khoản trước khi ghi. \\
\end{longtblr}

\section{Quy trình nghiệp vụ chính}

\subsection{Quy trình thêm sinh viên vào dự án}

Giảng viên thêm sinh viên bằng email. Server tìm tài khoản theo email, sau đó kiểm tra quyền quản lý dự án, trạng thái dự án, trạng thái tài khoản, vai trò sinh viên và tư cách thành viên hiện tại. Thao tác này được đồng bộ với các thay đổi tài khoản để tránh trường hợp một sinh viên vừa bị vô hiệu hóa ở yêu cầu khác nhưng vẫn được thêm vào dự án.

\figurelead{fig:add-student-workflow.png}{cụ thể hóa các bước kiểm tra cần thực hiện trước khi một sinh viên được ghi nhận là thành viên dự án}
\figAddStudentWorkflow

\subsection{Quy trình nộp bài và duyệt bài}

Bài nộp là dữ liệu có nhiều điều kiện đi kèm. Sinh viên phải đang hoạt động, là thành viên hiện tại của dự án, được phân công nhiệm vụ và dự án phải ở trạng thái active. Sau khi bài nộp được duyệt, tài nguyên và minh chứng gắn với bài nộp chuyển sang chỉ đọc để bảo toàn lịch sử đánh giá.

\figurelead{fig:review-workflow.png}{mô tả trình tự từ thời điểm sinh viên gửi báo cáo đến khi người duyệt ghi nhận quyết định cuối cùng}
\figReviewWorkflow

\subsection{Quy trình vòng đời dự án}

Vòng đời dự án là căn cứ quyết định các thao tác ghi. Ở mức phân tích yêu cầu, quy trình này được trình bày như một luồng nghiệp vụ từ thao tác của người quản lý tới phản hồi của hệ thống; phần kiểm soát kỹ thuật chi tiết được phân tích ở Chương 4 và Chương 5.

\figurelead{fig:project-lifecycle-workflow.png}{thể hiện các trạng thái dự án và ảnh hưởng của từng trạng thái tới thao tác nghiệp vụ được phép thực hiện}
\figProjectLifecycleWorkflow

\tableintro{tab:core-business-rules}{tổng hợp các luật nghiệp vụ bắt buộc trong những luồng dễ phát sinh sai lệch dữ liệu}
\begin{table}[H]
\caption{Luật nghiệp vụ bắt buộc cho các luồng chính}
\label{tab:core-business-rules}
\centering
\begin{tblr}{
  width=\textwidth,
  colspec={X[0.9,l]X[1.45,l]X[1.45,l]},
  rows={valign=t},
  row{1}={font=\bfseries,bg=black!7},
  hlines={0.35pt,black!24},
}
Luồng nghiệp vụ & Điểm dễ sai nếu làm thủ công & Quy tắc UniTrack áp dụng \\
Thêm thành viên & Sinh viên không còn hoạt động hoặc không thuộc dự án vẫn được giao việc & Server đọc trạng thái tài khoản mới nhất và chỉ nhận sinh viên active; khi xóa thành viên thì dọn liên kết phân công liên quan \\
Tạo nhiệm vụ & Nhiệm vụ không gắn mốc hoặc gán nhầm sinh viên ngoài dự án & Nhiệm vụ phải thuộc mốc cùng dự án; người được giao phải là thành viên active \\
Nộp bài & Một nhiệm vụ có nhiều bài chờ duyệt, hoặc sinh viên nộp sau khi dự án đã đóng & Ràng buộc một bài chờ duyệt và khóa vòng đời chặn request đã cũ \\
Duyệt bài & Hai giảng viên cùng duyệt một bài hoặc duyệt khi minh chứng chưa tải đủ & Server chỉ cho một quyết định cuối cùng; giao diện chặn lưu nếu dữ liệu hỗ trợ chưa sẵn sàng \\
Minh chứng & File minh chứng bị sửa/xóa sau khi đã là căn cứ đánh giá & API và trigger cơ sở dữ liệu chặn ghi sau khi duyệt, nhưng vẫn cho tải xuống \\
\end{tblr}
\end{table}

\section{Yêu cầu chức năng}

\tableintro{tab:functional-requirements}{tổng hợp các yêu cầu chức năng được suy ra từ khảo sát, tác nhân và use case}
\begin{longtblr}[
  caption={Danh sách yêu cầu chức năng của UniTrack},
  label={tab:functional-requirements},
]{
  width=\textwidth,
  colspec={X[0.45,l]X[0.95,l]X[2.15,l]},
  rowhead=1,
  rows={valign=t},
  row{1}={font=\bfseries,bg=black!7},
  hlines={0.35pt,black!24},
}
Mã & Nhóm & Nội dung yêu cầu \\
FR01 & Xác thực & Người dùng đăng nhập, đăng xuất và xem current user; session cookie được tạo, thu hồi và kiểm tra trạng thái tài khoản. \\
FR02 & Quản lý tài khoản & Quản trị viên tạo, tìm kiếm, cập nhật vai trò/trạng thái, đặt mật khẩu và xử lý các tài khoản đang liên quan tới công việc mở. \\
FR03 & Dashboard & Hệ thống hiển thị danh sách bài chờ duyệt, nhiệm vụ quá hạn, dự án cần chú ý hoặc nhiệm vụ cần xử lý theo từng vai trò. \\
FR04 & Workspace & Giảng viên/quản trị viên quản lý thư mục và danh sách dự án; sinh viên chỉ thấy dự án mình tham gia. \\
FR05 & Dự án & Giảng viên/quản trị viên tạo, sửa dự án, cập nhật trạng thái vòng đời, gán thư mục và xem tóm tắt tiến độ. \\
FR06 & Thành viên & Người quản lý thêm sinh viên active bằng email, đổi vai trò leader/member, xóa thành viên và dọn các phân công liên quan. \\
FR07 & Mốc kế hoạch & Người quản lý tạo, sửa, xóa khi hợp lệ và sắp xếp lại mốc kế hoạch trong phạm vi dự án. \\
FR08 & Nhiệm vụ & Người quản lý tạo/sửa nhiệm vụ thuộc mốc kế hoạch, đặt hạn nộp, mức ưu tiên và phân công sinh viên active trong dự án. \\
FR09 & Bài nộp & Sinh viên được phân công nộp báo cáo tiến độ khi dự án active và chưa có bài chờ duyệt trùng nhiệm vụ. \\
FR10 & Duyệt bài & Người quản lý duyệt một bài đang chờ duyệt một lần duy nhất, đồng bộ trạng thái bài nộp và trạng thái tiến độ nhiệm vụ. \\
FR11 & Tài nguyên & Người có quyền xem dự án được thêm đường dẫn tài nguyên vào dự án, mốc kế hoạch, nhiệm vụ hoặc bài nộp đang chờ duyệt theo trạng thái vòng đời. \\
FR12 & Minh chứng & Người nộp hoặc người quản lý được tải lên, tải xuống, xóa minh chứng cho bài đang chờ duyệt; minh chứng đã duyệt chỉ còn đọc/tải xuống. \\
FR13 & Phân quyền & Giao diện và server kiểm tra quyền theo vai trò, quan hệ với dự án, trạng thái tài khoản và vòng đời dự án. \\
FR14 & Toàn vẹn dữ liệu & Cơ sở dữ liệu bảo vệ quan hệ cùng dự án giữa nhiệm vụ, bài nộp và minh chứng; bảo đảm một bài chờ duyệt và không sửa minh chứng sau khi duyệt. \\
FR15 & Kiểm thử & Có kiểm thử backend, kiểm thử trình duyệt bằng Playwright, kiểm tra lint/build giao diện và kiểm tra migration cơ sở dữ liệu. \\
\end{longtblr}

\section{Đặc tả use case tiêu biểu}

Phần đặc tả chi tiết chỉ chọn bảy use case trọng tâm. Các thao tác còn lại được mô tả ở mức yêu cầu chức năng để báo cáo giữ đúng trọng tâm nghiệp vụ.

\subsection{UC01 -- Đăng nhập hệ thống}

\tableintro{tab:uc-login}{đặc tả tiền điều kiện, luồng chính, luồng thay thế và hậu điều kiện của use case đăng nhập}
\begin{longtblr}[
  caption={Đặc tả use case đăng nhập hệ thống},
  label={tab:uc-login},
]{
  width=\textwidth,
  colspec={X[0.8,l]X[2.45,l]},
  rowhead=1,
  rows={valign=t},
  row{1}={font=\bfseries,bg=black!7},
  hlines={0.35pt,black!24},
}
Thuộc tính & Mô tả \\
Tác nhân & Quản trị viên, giảng viên hoặc sinh viên đã có tài khoản active \\
Tiền điều kiện & API hoạt động, tài khoản tồn tại, mật khẩu đúng, trạng thái tài khoản là \code{active} \\
Luồng chính & Người dùng nhập email và mật khẩu; hệ thống kiểm tra thông tin đăng nhập, trạng thái tài khoản và giới hạn đăng nhập; nếu hợp lệ, hệ thống tạo phiên đăng nhập và chuyển người dùng tới dashboard theo vai trò. \\
Luồng thay thế & Sai thông tin đăng nhập, tài khoản không hoạt động, dữ liệu form không hợp lệ hoặc vượt giới hạn đăng nhập sẽ bị từ chối và giao diện hiển thị thông báo lỗi phù hợp. \\
Hậu điều kiện & Người dùng có phiên đăng nhập hợp lệ; giao diện biết vai trò hiện tại và chuyển tới dashboard hoặc trang đang chờ. \\
\end{longtblr}

\subsection{UC02 -- Tạo dự án}

\tableintro{tab:uc-create-project}{đặc tả use case tạo dự án với các điều kiện liên quan đến quyền và thư mục Workspace}
\begin{longtblr}[
  caption={Đặc tả use case tạo dự án},
  label={tab:uc-create-project},
]{
  width=\textwidth,
  colspec={X[0.8,l]X[2.45,l]},
  rowhead=1,
  rows={valign=t},
  row{1}={font=\bfseries,bg=black!7},
  hlines={0.35pt,black!24},
}
Thuộc tính & Mô tả \\
Tác nhân & Giảng viên hoặc quản trị viên \\
Tiền điều kiện & Người dùng đã đăng nhập; nếu chọn thư mục thì thư mục active và chủ thư mục khớp giảng viên phụ trách dự án \\
Luồng chính & Người dùng mở hộp thoại tạo dự án, nhập tên, chủ đề, mô tả, khoảng thời gian và thư mục nếu có. Hệ thống kiểm tra quyền tạo dự án, thông tin thời gian và việc gán thư mục, sau đó tạo dự án mới. \\
Luồng thay thế & Tên rỗng, ngày không hợp lệ, thư mục không phù hợp hoặc người dùng không đủ quyền sẽ bị từ chối; giao diện giữ lại dữ liệu để người dùng sửa. \\
Hậu điều kiện & Dự án mới xuất hiện trong workspace, trang chi tiết thư mục hoặc danh sách dự án độc lập và được đưa vào các thống kê liên quan. \\
\end{longtblr}

\subsection{UC03 -- Tạo nhiệm vụ chính thức}

\tableintro{tab:uc-create-assignment}{đặc tả use case tạo nhiệm vụ chính thức trong phạm vi mốc kế hoạch và thành viên dự án}
\begin{longtblr}[
  caption={Đặc tả use case tạo nhiệm vụ chính thức},
  label={tab:uc-create-assignment},
]{
  width=\textwidth,
  colspec={X[0.8,l]X[2.45,l]},
  rowhead=1,
  rows={valign=t},
  row{1}={font=\bfseries,bg=black!7},
  hlines={0.35pt,black!24},
}
Thuộc tính & Mô tả \\
Tác nhân & Giảng viên hoặc quản trị viên có quyền quản lý dự án \\
Tiền điều kiện & Dự án active; mốc kế hoạch thuộc cùng dự án; sinh viên được giao là thành viên active của dự án \\
Luồng chính & Người quản lý mở chức năng tạo nhiệm vụ trong kế hoạch, chọn mốc, nhập tiêu đề, mô tả, hạn nộp, mức ưu tiên và danh sách sinh viên. Hệ thống kiểm tra dự án còn cho phép giao việc, mốc thuộc dự án và sinh viên được giao là thành viên hợp lệ, sau đó tạo nhiệm vụ. \\
Luồng thay thế & Dự án không còn cho phép giao việc, mốc không cùng dự án hoặc sinh viên được giao không hợp lệ sẽ bị từ chối; giao diện tải lại dữ liệu khi trạng thái dự án đã thay đổi. \\
Hậu điều kiện & Nhiệm vụ xuất hiện dưới mốc kế hoạch, sinh viên được giao nhìn thấy trong dashboard và thống kê dự án được cập nhật. \\
\end{longtblr}

\subsection{UC04 -- Nộp báo cáo và minh chứng}

\tableintro{tab:uc-submit-progress}{đặc tả use case sinh viên nộp báo cáo tiến độ và gắn minh chứng cho nhiệm vụ được phân công}
\begin{longtblr}[
  caption={Đặc tả use case nộp báo cáo và minh chứng},
  label={tab:uc-submit-progress},
]{
  width=\textwidth,
  colspec={X[0.8,l]X[2.45,l]},
  rowhead=1,
  rows={valign=t},
  row{1}={font=\bfseries,bg=black!7},
  hlines={0.35pt,black!24},
}
Thuộc tính & Mô tả \\
Tác nhân & Sinh viên được phân công nhiệm vụ \\
Tiền điều kiện & Sinh viên active, là thành viên dự án, được phân công vào nhiệm vụ, dự án active, nhiệm vụ chưa hoàn tất, chưa có bài chờ duyệt cùng nhiệm vụ và file minh chứng không vượt 10 MB. \\
Luồng chính & Sinh viên mở trang nhiệm vụ, nhập tiêu đề, mô tả và vướng mắc rồi gửi bài nộp. Hệ thống kiểm tra sinh viên còn thuộc dự án, còn được phân công và nhiệm vụ chưa có bài đang chờ duyệt, sau đó ghi nhận bài nộp. Sinh viên có thể gắn file minh chứng cho bài đang chờ duyệt. \\
Luồng thay thế & Dự án đã đóng ghi, sinh viên bị xóa khỏi nhóm, nhiệm vụ đã hoàn tất, có bài chờ duyệt trùng, file quá lớn hoặc bài đã được duyệt sẽ bị từ chối; giao diện tải lại dữ liệu nhiệm vụ/dự án. \\
Hậu điều kiện & Bài nộp xuất hiện trong lịch sử và danh sách chờ duyệt; minh chứng gắn với bài nộp đang chờ duyệt và sẵn sàng cho người duyệt kiểm tra. \\
\end{longtblr}

\subsection{UC05 -- Duyệt bài nộp}

\tableintro{tab:uc-review-submission}{đặc tả use case giảng viên hoặc quản trị viên duyệt bài nộp đang chờ xử lý}
\begin{longtblr}[
  caption={Đặc tả use case duyệt bài nộp},
  label={tab:uc-review-submission},
]{
  width=\textwidth,
  colspec={X[0.8,l]X[2.45,l]},
  rowhead=1,
  rows={valign=t},
  row{1}={font=\bfseries,bg=black!7},
  hlines={0.35pt,black!24},
}
Thuộc tính & Mô tả \\
Tác nhân & Giảng viên hoặc quản trị viên có quyền quản lý dự án \\
Tiền điều kiện & Bài nộp đang chờ duyệt, dự án chưa archived, người duyệt còn quyền quản lý \\
Luồng chính & Người duyệt mở trang duyệt bài, xem mô tả nhiệm vụ, ngữ cảnh, tài nguyên và minh chứng, chọn chấp nhận/yêu cầu sửa/từ chối, nhập hướng dẫn khi cần và gửi quyết định. Hệ thống ghi nhận quyết định cuối cùng và cập nhật trạng thái nhiệm vụ. \\
Luồng thay thế & Bài nộp đã được người khác duyệt, quyết định yêu cầu sửa/từ chối nhưng thiếu hướng dẫn, hoặc dữ liệu minh chứng chưa sẵn sàng sẽ bị chặn. \\
Hậu điều kiện & Bài nộp có quyết định duyệt cuối cùng; tài nguyên và minh chứng của bài nộp chuyển sang chỉ đọc nhưng vẫn tải xuống được. \\
\end{longtblr}

\subsection{UC06 -- Quản lý tài khoản}

\tableintro{tab:uc-admin-accounts}{đặc tả use case quản trị tài khoản và các điều kiện giữ an toàn session}
\begin{longtblr}[
  caption={Đặc tả use case quản lý tài khoản},
  label={tab:uc-admin-accounts},
]{
  width=\textwidth,
  colspec={X[0.8,l]X[2.45,l]},
  rowhead=1,
  rows={valign=t},
  row{1}={font=\bfseries,bg=black!7},
  hlines={0.35pt,black!24},
}
Thuộc tính & Mô tả \\
Tác nhân & Quản trị viên \\
Tiền điều kiện & Quản trị viên đã đăng nhập; hệ thống còn ít nhất một quản trị viên active sau thao tác. \\
Luồng chính & Quản trị viên mở trang quản lý tài khoản, tìm kiếm hoặc tạo tài khoản mới, chọn vai trò/trạng thái và đặt mật khẩu khi cần. Hệ thống kiểm tra quyền quản trị, tránh làm mất quản trị viên hoạt động cuối cùng, thu hồi phiên đăng nhập khi đổi mật khẩu/vô hiệu hóa và ghi nhật ký thao tác. \\
Luồng thay thế & Email trùng, vai trò không hợp lệ, cố vô hiệu hóa quản trị viên active cuối cùng hoặc tài khoản đang phụ trách công việc mở mà chưa chọn phương án chuyển giao sẽ bị từ chối. \\
Hậu điều kiện & Danh sách tài khoản được cập nhật; session liên quan bị thu hồi khi cần; dashboard và cache tài khoản được làm mới. \\
\end{longtblr}

\subsection{UC07 -- Thêm sinh viên vào dự án}

\tableintro{tab:uc-add-student}{đặc tả use case thêm sinh viên active vào dự án với kiểm tra trạng thái tài khoản mới nhất}
\begin{longtblr}[
  caption={Đặc tả use case thêm sinh viên vào dự án},
  label={tab:uc-add-student},
]{
  width=\textwidth,
  colspec={X[0.8,l]X[2.45,l]},
  rowhead=1,
  rows={valign=t},
  row{1}={font=\bfseries,bg=black!7},
  hlines={0.35pt,black!24},
}
Thuộc tính & Mô tả \\
Tác nhân & Giảng viên hoặc quản trị viên có quyền quản lý dự án \\
Tiền điều kiện & Dự án active hoặc on hold theo quy tắc nhóm; tài khoản sinh viên đã tồn tại và đang active. \\
Luồng chính & Người quản lý mở popover nhóm, nhập email sinh viên và gửi yêu cầu thêm. Hệ thống kiểm tra quyền quản lý, trạng thái dự án, vai trò và trạng thái tài khoản sinh viên rồi thêm sinh viên nếu chưa là thành viên. \\
Luồng thay thế & Email không tồn tại, tài khoản không phải sinh viên, tài khoản inactive, dự án archived/completed hoặc sinh viên đã là thành viên sẽ trả lỗi và giao diện giữ nguyên danh sách hiện tại. \\
Hậu điều kiện & Sinh viên thấy dự án trong Workspace; người quản lý có thể phân công sinh viên vào nhiệm vụ sau đó. \\
\end{longtblr}

\section{Yêu cầu phi chức năng}

\tableintro{tab:nonfunctional-requirements}{tổng hợp các yêu cầu phi chức năng quan trọng về sử dụng, bảo mật, toàn vẹn dữ liệu, bảo trì, kiểm thử và triển khai}
\begin{longtblr}[
  caption={Yêu cầu phi chức năng của UniTrack},
  label={tab:nonfunctional-requirements},
]{
  width=\textwidth,
  colspec={X[0.85,l]X[2.55,l]},
  rowhead=1,
  rows={valign=t},
  row{1}={font=\bfseries,bg=black!7},
  hlines={0.35pt,black!24},
}
Nhóm yêu cầu & Mô tả \\
Tính dễ dùng & Giao diện phải dẫn người dùng tới việc cần làm tiếp theo: giảng viên/quản trị viên thấy bài chờ duyệt và nhiệm vụ quá hạn; sinh viên thấy nhiệm vụ cần xử lý. \\
Bảo mật session & Session cookie phải được cấu hình an toàn theo môi trường; các request thay đổi dữ liệu phải đến từ nguồn tin cậy; hệ thống phải thu hồi session khi tài khoản thay đổi. \\
Phân quyền & Server là nguồn quyết định cuối cùng cho vai trò, trạng thái tài khoản, quan hệ với dự án và vòng đời dự án. \\
Toàn vẹn dữ liệu & PostgreSQL phải có khóa ngoại, chỉ mục duy nhất và trigger để giữ nhiệm vụ, bài nộp, tài nguyên/minh chứng cùng dự án và không bị sửa sau khi duyệt. \\
Khả năng bảo trì & Mã nguồn tổ chức theo lát tính năng; phần mô tả kỹ thuật phải bám sát triển khai hiện tại. \\
Khả năng kiểm thử & Kiểm thử phía server phải bao phủ các ràng buộc nghiệp vụ; Playwright kiểm tra những màn hình rủi ro cao từ góc nhìn người dùng. \\
Triển khai & Hệ thống chạy cục bộ với PostgreSQL và có hướng triển khai Vercel, Render, Neon, Cloudflare R2; môi trường production phải có cấu hình đăng nhập, CORS, storage và tài khoản quản trị ban đầu an toàn. \\
\end{longtblr}

Các yêu cầu chính của UniTrack đã được xác định theo ba vai trò, các use case trọng tâm và các ràng buộc nghiệp vụ cần bảo vệ. Những yêu cầu này là cơ sở để Chương 3 lựa chọn công nghệ phù hợp và Chương 4 thiết kế chi tiết kiến trúc, dữ liệu, giao diện và kiểm thử.
